\documentclass[11pt,a4paper]{article}
\usepackage{fontspec}\setmainfont{Arial}\setsansfont{Arial}\setmonofont{Arial}
\usepackage{polyglossia}\setdefaultlanguage{russian}\setotherlanguage{english}
\usepackage{geometry,amsmath,amssymb,booktabs,longtable,tabularx,array,enumitem,xcolor,hyperref,fancyhdr,tikz}
\usetikzlibrary{arrows.meta,positioning,shapes.geometric}
\geometry{margin=18mm}\setlength{\headheight}{14pt}\setlength{\parindent}{0pt}\setlength{\parskip}{6pt}
\setlist{leftmargin=*,nosep}
\definecolor{navy}{HTML}{17365D}\definecolor{teal}{HTML}{167D8D}\definecolor{pale}{HTML}{EEF5F7}\definecolor{orange}{HTML}{D87800}\definecolor{green}{HTML}{287A4A}
\hypersetup{colorlinks=true,linkcolor=navy,urlcolor=teal}
\pagestyle{fancy}\fancyhf{}\lhead{MATEMAFIJA}\rhead{Объяснение модели}\cfoot{\thepage}
\newcommand{\code}[1]{\texttt{#1}}
\newcommand{\explain}[1]{\vspace{1mm}\noindent\fcolorbox{teal}{pale}{\parbox{0.94\linewidth}{#1}}\vspace{1mm}}
\title{\vspace{-1cm}\textbf{MATEMAFIJA}\\[2mm]\Large Что это, зачем это и как оно работает\\[2mm]\large Дополнительное объяснение transition и response model}
\author{Команда MATEMAFIJA}
\date{9 октября 2026}

\begin{document}\maketitle
\explain{\textbf{Короткий ответ на главный вопрос.} $h_1^{release}$ и $h_1^{start}$ --- не наблюдаемые величины и не новая цена. Это две \textbf{условные интенсивности переходов} после применения гипотетической реакции на adaptive price. В текущем коде они преобразуются в вероятности, затем в stochastic release/start events, затем в projected EVSE counts и scarcity metrics. Они влияют только на assumption-driven simulation. Они не изменяют SQLite, observed history, F1 или сам OSA signal.}
\tableofcontents\newpage

\section{Почему этот слой вообще нужен}
Pricing policy отвечает только на вопрос: \emph{какой ценовой сигнал показать?} Но город интересует другой вопрос: \emph{изменится ли доступность зарядок?} Между ценой и доступностью отсутствует автоматически известная связь.

Поэтому архитектура разделена:
\[
\text{observed state}\rightarrow\text{price decision}\rightarrow
\boxed{\text{behavioural response assumption}}\rightarrow
\text{simulated state}\rightarrow\text{scarcity metrics}.
\]
Без центрального response block формула OSA могла бы показать только цену. Она не могла бы освободить ни одного EVSE даже внутри модели.

\section{Сначала различим четыре разных объекта}
\begin{tabularx}{\textwidth}{p{3cm}X p{3cm}}
\toprule Объект & Что означает & Provenance \\\midrule
$h_0(t)$ & baseline transition hazard, оценённый из qualified historical transitions & estimated \\
$d_i(t)$ & применимый policy adjustment в EUR/min после $\lambda$ и rollout & policy-defined \\
$\theta_r,\theta_a,\tau$ & предполагаемая чувствительность и задержка & assumed \\
$h_1(t),q_1(t)$ & изменённые hazard/probability при данном сценарии & simulated \\
\bottomrule
\end{tabularx}

Ключевой смысл: $h_0$ описывает «как система обычно переходила», $d$ --- «какой сигнал создаёт policy», $\theta$ --- «как мы допускаем, что люди реагируют», а $h_1$ --- «каким станет transition process, если это допущение верно».

\section{Что такое hazard}
Hazard $h$ --- не вероятность и не число машин. Это мгновенная условная интенсивность события с единицей $1/\text{hour}$. Для шага длительностью $\Delta t$ она превращается в вероятность:
\[
q=1-\exp(-h\Delta t).
\]
При маленьком шаге $q\approx h\Delta t$, но код использует точную формулу. Например, при $h=0.6/\text{hour}$ и $\Delta t=5$ min $=1/12$ hour:
\[
q=1-e^{-0.6/12}\approx0.0488.
\]
То есть каждый подходящий EVSE имеет примерно 4.88\% вероятности такого перехода на следующем пятиминутном шаге. Это не означает, что событие обязательно произойдёт.

\section{Два канала: release и accepted start}
\subsection{Release channel}
\[
\boxed{h_1^{release}(t)=h_0^{release}(t)\exp\{\theta_r d_i(t-\tau)\}}
\]
\textbf{Что это:} условная интенсивность перехода \code{charging $\to$ available} после price response.\\
\textbf{Зачем:} выразить гипотезу, что положительный price adjustment ускоряет освобождение EVSE.\\
\textbf{Почему знак плюс:} при $\theta_r,d>0$ множитель $e^{\theta_r d}>1$.\\
\textbf{Что это не доказывает:} окончание энергии, post-charge idle и мотив конкретного водителя не наблюдаются. Поэтому название --- \emph{release proxy}.

\subsection{Accepted-start channel}
\[
\boxed{h_1^{start}(t)=h_0^{start}(t)\exp\{-\theta_a d_i(t-\tau)\}}
\]
\textbf{Что это:} условная интенсивность перехода \code{available $\to$ charging}.\\
\textbf{Зачем:} учесть, что более высокая цена может уменьшить число принятых данной станцией charging starts.\\
\textbf{Почему знак минус:} при $\theta_a,d>0$ множитель $e^{-\theta_a d}<1$.\\
\textbf{Что это не означает:} мы не видим всех arrivals. Это accepted-start proxy; отвергнутая/отложенная попытка могла исчезнуть, перенестись или произойти позднее.

\explain{Release улучшает availability через большее число переходов charging $\to$ available. Start response также может улучшить availability, но ценой меньшего accepted use. Поэтому проект обязан показывать \code{deferred\_or\_displaced\_start\_proxy}; иначе можно было бы ошибочно назвать подавление спроса эффективностью.}

\section{Что такое выделенный $D_i(t)$}
$D_i(t)$ --- время с начала \textbf{непрерывно наблюдаемого qualifying scarcity condition}. Это reconstructed lower bound, а не время простоя автомобиля.

Condition продолжается, только если одновременно: snapshot пригоден; scarcity выше trigger; есть хотя бы один charging EVSE; equipment loss не превышает guard; нет collection gap больше допустимого. При нарушении duration сбрасывается.

\textbf{Зачем:} не повышать цену из-за единственного шумного snapshot. Устойчивый дефицит сильнее кратковременного.

\section{Что делает $P_i(t)^\gamma R_i(t)$}
\[
P_i=\operatorname{clip}\left(\frac{C_i-C_{trigger}}{1-C_{trigger}},0,1\right),\qquad
R_i=\operatorname{clip}\left(\frac{D_i-T_{grace}}{T_{ramp}},0,1\right).
\]
$P_i$ отвечает «насколько сильный scarcity», $R_i$ --- «насколько долго он длится после grace». $\gamma=1.2$ немного подавляет слабое pressure и подчёркивает сильное. Их произведение
\[
S_i=P_i^\gamma R_i
\]
лежит в $[0,1]$. Затем $\Delta p_{raw}=0.05S_i$ EUR/min. Следовательно, cap математически гарантирован.

Пример: $C=0.75$, trigger $0.5$, $D=30$ min. Тогда $P=0.5$, $R=(30-15)/30=0.5$, $S\approx0.5^{1.2}\cdot0.5\approx0.218$, а raw adjustment около 0.0109 EUR/min.

\section{Что делает выделенный $d_i(t)$}
\[
d_i(t)=R_i^{rollout}(t)\lambda\Delta p_{raw}(i,t).
\]
Это \textbf{финальный response driver}, уже пригодный для behavioural layer.
\begin{itemize}
\item $R_i^{rollout}=0/1$: policy ещё не/уже развёрнута на site.
\item $\lambda\in[0,1]$: intensity. При $\lambda=0$ adaptive effect отсутствует; при 1 применяется весь raw adjustment.
\item $\Delta p_{raw}$: output OSA-v1.
\end{itemize}
Coverage $\kappa$ выбирает долю sites/EVSE, rollout определяет где и когда, а $\lambda$ --- силу. Их нельзя смешивать.

\section{Как $h_1$ влияет в текущем коде: цепочка по шагам}
\begin{center}
\begin{tikzpicture}[node distance=4mm,every node/.style={font=\small},b/.style={draw,rounded corners,fill=pale,text width=13cm,align=center,minimum height=8mm},a/.style={-{Stealth},thick,teal}]
\node[b](x){OSA + F1 + adoption $\rightarrow d_i(t)$};
\node[b,below=of x](h){$h_0$ из transition model; response formula $\rightarrow h_1^{release},h_1^{start}$};
\node[b,below=of h](q){$q_0=1-e^{-h_0\Delta t}$ и $q_1=1-e^{-h_1\Delta t}$};
\node[b,below=of q](u){Для каждого source EVSE генерируется $u\sim U(0,1)$; то же $u$ используется в fixed и adaptive};
\node[b,below=of u](n){Событие происходит, если $u<q$: получаем release/start counts};
\node[b,below=of n](s){$A'=A+releases-starts$, $C'=C-releases+starts$};
\node[b,below=of s](m){Одинаковые scarcity metrics $\rightarrow\Delta M\rightarrow$ robustness};
\foreach \x/\y in {x/h,h/q,q/u,u/n,n/s,s/m}\draw[a](\x)--(\y);
\end{tikzpicture}
\end{center}

То есть $h_1$ влияет не напрямую на metric, а через четыре преобразования: hazard $\to$ probability $\to$ random transition count $\to$ projected state $\to$ metric.

\section{Числовой пример влияния}
Пусть шаг 5 минут, $d=0.03$ EUR/min.

\textbf{Release:} $h_0^r=0.6$/h, $\theta_r=20$ min/EUR.
\[
e^{\theta_r d}=e^{0.6}=1.822,\quad h_1^r=1.093/\text{h}.
\]
\[
q_0^r=1-e^{-0.6/12}=0.0488,\quad q_1^r=1-e^{-1.093/12}=0.0871.
\]
Для 10 charging EVSE expected releases на шаге меняются примерно с 0.49 до 0.87.

\textbf{Start:} $h_0^a=0.4$/h, $\theta_a=15$ min/EUR.
\[
e^{-\theta_a d}=e^{-0.45}=0.638,\quad h_1^a=0.255/\text{h};
\]
\[
q_0^a\approx0.0328,\quad q_1^a\approx0.0210.
\]
Для 10 available EVSE expected accepted starts меняются примерно с 0.33 до 0.21.

В одной replication целые counts могут совпасть из-за случайности. Поэтому выполняется много paired replications и оценивается частота улучшения.

\section{Почему common random numbers важны}
Если fixed и adaptive получили бы независимые random draws, разница могла бы возникнуть просто из-за удачи. Сейчас один и тот же $u$ сравнивается с $q_0$ и $q_1$.

Например, $q_0^r=0.049$, $q_1^r=0.087$:
\begin{itemize}
\item $u=0.03$: release в обоих arms;
\item $u=0.07$: только adaptive release --- эффект response assumption;
\item $u=0.20$: release ни в одном arm.
\end{itemize}
Так paired difference лучше изолирует влияние policy/response parameters.

\section{Что означает one-step projection}
Для каждого baseline snapshot строится независимый следующий условный state. Projected state из момента $t$ \textbf{не становится} source state в $t+1$.

\textbf{Зачем так сделано сейчас:} это безопасный MVP, который не накапливает ошибки и не требует пока полной demand/queue/spatial dynamics.\\
\textbf{Цена ограничения:} нельзя интерпретировать соседние projected snapshots как связную траекторию; поэтому simulated episode duration отключена. Можно оценивать snapshot-level availability/scarcity, но не длительность simulated shortage episode.

\section{Что такое grid, replications и $\Delta M$}
Grid перебирает все комбинации:
\[
\lambda\times\theta_r\times\theta_a\times\tau\times\text{replication}.
\]
Одна комбинация без replication --- scenario assumptions. Replications отвечают за stochastic uncertainty. Для каждого запуска:
\[
\Delta M=M_{adaptive}-M_{fixed}.
\]
Знак зависит от metric: для zero-availability/scarcity отрицательный --- лучше; для availability положительный --- лучше.

\section{Wilson bound, robustness и minimum threshold}
Если adaptive лучше в $k$ из $n$ replications, $\hat p=k/n$. Но просто $\hat p\ge0.8$ при малом $n$ слишком оптимистично. Wilson interval учитывает uncertainty доли.

Текущий robust criterion: минимум 20 valid replications и \textbf{нижняя 95\% Wilson bound $\ge0.80$}. Это означает: даже консервативная граница должна пройти 80\%.

\textbf{Minimum response threshold} --- наименьшее протестированное $\theta_r$, прошедшее criterion при фиксированных $\lambda,\theta_a,\tau$. Это не истинная elasticity и не аналитически найденный optimum; это threshold на выбранном grid.

\section{Что происходит прямо сейчас, а что пока нет}
\begin{tabularx}{\textwidth}{X X}
\toprule Уже реализовано в коде & Пока не является эмпирическим результатом \\\midrule
формулы release/start hazard shifts & значения $\theta$ не оценены по Vilnius intervention \\
hazard-to-probability conversion & нет доказанной causal response \\
paired Bernoulli draws & нет observed queue/waiting/idle \\
projected counts и scarcity metrics & нет stateful multi-step trajectory \\
grid, replications, Wilson robustness & current data window ещё не проходит 7-day gate \\
withholding при неготовых inputs & только illustrative response preview до readiness \\
\bottomrule
\end{tabularx}

Главное практическое следствие: $h_1$ \textbf{работает как программный механизм}, но его численные outputs пока не должны подаваться как результат о реальном поведении водителей Вильнюса. UI отделяет labelled illustrative mechanism preview от data-qualified effectiveness result.

\section{Где это находится в репозитории}
\begin{longtable}{p{5.5cm}p{10.5cm}}
\toprule Файл & Ответственность \\\midrule\endhead
\code{src/transition\_model.py} & оценивает qualified $h_0^{release}$ и $h_0^{start}$.\\
\code{src/adaptive\_policy.py} & рассчитывает OSA raw adjustment из pressure/persistence.\\
\code{src/policy\_integration.py} & применяет $\lambda$, rollout и создаёт $d_i(t)$.\\
\code{src/release\_response.py} & рассчитывает release multiplier, $h_1^r$, $q_0^r$, $q_1^r$.\\
\code{src/arrival\_response.py} & рассчитывает start multiplier, $h_1^a$, $q_0^a$, $q_1^a$.\\
\code{src/simulation\_grid.py} & draws, counts, one-step states, paired runs.\\
\code{src/simulation\_evaluation.py} & одинаковые scarcity metrics и $\Delta M$.\\
\code{src/robustness.py} & Wilson bounds, robust region, tested threshold.\\
\bottomrule
\end{longtable}

\section{Как это объяснить жюри за 30 секунд}
«Цена сама по себе не изменяет доступность. Поэтому мы отделяем pricing mechanism от response assumption. Из истории оцениваем baseline rates переходов charging-to-available и available-to-charging. Затем явно варьируем неизвестную sensitivity: положительная price adjustment может ускорить release и уменьшить accepted starts. Эти hazards превращаются в paired stochastic transitions и проходят через одни scarcity metrics. Мы не выдаём assumptions за поведение Vilnius; мы ищем threshold, который затем можно проверить пилотом.»

\section{Пять проверочных вопросов}
\begin{enumerate}
\item Если $d=0$ или $\theta=0$, почему $h_1=h_0$?
\item Почему больше release и меньше starts оба улучшают availability, но только второе создаёт access trade-off?
\item Почему $h$ нельзя назвать вероятностью без $\Delta t$?
\item Почему one-step model не позволяет считать simulated episode duration?
\item Почему robust threshold остаётся conditional, даже после тысячи replications?
\end{enumerate}
Если ответы ясны, центральная логика simulator понятна.

\vfill\small\textit{Документ объясняет реализацию репозитория на 9 октября 2026. Все behavioural sensitivities остаются assumed до пилотной или иной надёжной валидации.}
\end{document}
